% Desarrollo reunido K31-07--K31-15. Inserción anterior a registro_k y alpha.
% Mantiene explícita la interfaz del repertorio terminal de ocho bloques.
% Propietarios y localizadores: metadata/K_ALPHA_PROPIETARIOS.md.
\section{Regional determination of the dodecaphasic descriptor and selection of the coupling register}
\label{chap:despliegue-selector-k}
\chaptermark{Dodecaphase coupling register}

The determination of the coupling register articulates two constructions that should be developed in their effective order. The first produces the selector's antecedents from the regional expressions of the TPK: six-trit bands, column signatures, conversion calendar, oriented descriptor and functional panel. The second operates on these antecedents and on the terminal repertoire of eight blocks, enumerates their admissible orderings and selects a unique register through two simultaneous conditions: exceptional incidence and reading heterotypy. The register's value appears at the end of this composition.

The present development presupposes the regional expressions obtained in the preceding chapters from APP--TRIT--TPK. Its purpose is to make explicit every intermediate operation leading from those expressions to the finite selector. The realization of the same register through signed incidences, Hadamard inversion and expression of the fine-structure constant are then incorporated. The three constructions preserve their respective domains: register selection, recovery from its channels and arithmetic evaluation of its expressions.

\subsection{Regional bands, encoding, and integer signatures}

\subsubsection{Domains and encoding of the regional bands}

The channel order is fixed as
\[
 \chi\in\{\pi,e,\varphi\},
 \qquad B_\chi(t)\in\mathbb F_3^6,
 \qquad t\in\mathbb N.
\]
The channel symbols identify the regions already generated. Each $B_\chi(t)$ preserves its position in its sequence of expressions; index $t$ counts ternary blocks. The natural evaluation of a band $b=(b_0,\ldots,b_5)$ is
\begin{equation}
 \operatorname{enc}_6(b)=
 243b_0+81b_1+27b_2+9b_3+3b_4+b_5,
 \qquad 0\leq\operatorname{enc}_6(b)<729.
 \label{eq:kdes-enc6}
\end{equation}
Its fixed-width inverse is
\[
 \operatorname{dig}_j(n)=
 \left\lfloor\frac{n}{3^{5-j}}\right\rfloor\bmod3,
 \qquad 0\leq j<6.
\]
Leading zeros are thus retained: the words $002001$ and $2001$
could share an abbreviated integer notation, but the type of the
former contains six positions. All subsequent concatenations
use this width.

In the window of one hundred blocks, the six-hundred-trit expression of channel $\chi$ determines an integer $P_\chi^{(600)}$. Extraction is performed through
\begin{equation}
 b_\chi(t)=
 \left\lfloor\frac{P_\chi^{(600)}}{729^{99-t}}\right\rfloor\bmod729,
 \qquad 0\leq t<100,
 \qquad B_\chi(t)=\operatorname{enc}_6^{-1}(b_\chi(t)).
 \label{eq:kdes-extraccion}
\end{equation}
The six-hundred-trit integer is an expression of the regional producer; the hundred-row table is the evaluation of \eqref{eq:kdes-extraccion}. This separation allows a historical table to be replaced by its producer without modifying the selector that consumes it.

\subsubsection{Residues, quotients, occupancy, and orientation}

For each column $j$ and time $t$, introduce
\begin{align}
 T_j(t)&=B_\pi(t)_j+B_e(t)_j+B_\varphi(t)_j,\\
 q_j(t)&=T_j(t)\bmod3,
 &c_j(t)&=\left\lfloor T_j(t)/3\right\rfloor,\\
 a_j(t)&=\bigl(B_\pi(t)_j+B_e(t)_j-B_\varphi(t)_j\bigr)\bmod3,\\
 w_j(t)&=\mathbf1_{B_\pi(t)_j\ne0}
       +\mathbf1_{B_e(t)_j\ne0}
       +\mathbf1_{B_\varphi(t)_j\ne0},
 &\bar w_j(t)&=w_j(t)\bmod3.
 \label{eq:kdes-firmas}
\end{align}
The subtraction in the third expression is performed in the integers and
then reduced. In an implementation over natural numbers, the equivalent expression
is $(B_\pi+B_e+3-B_\varphi)\bmod3$; truncated subtraction would change the reader.

\begin{proposition}[Reconstruction of the column sum]
For all admissible bands,
\[
 T_j(t)=q_j(t)+3c_j(t),\qquad
 0\leq q_j(t)<3,\quad 0\leq c_j(t)\leq2,
 \quad0\leq w_j(t)\leq3.
\]
\end{proposition}
\begin{proof}
Each column sums three representatives from $\{0,1,2\}$, so that
$0\leq T_j(t)\leq6$. Euclidean division by three yields the stated residue
and quotient. Occupancy counts three binary conditions;
its bound follows directly from that definition.
\end{proof}

This identity preserves the quotient of the local sum. Reconstructing
the three individual bands requires the bands themselves or an additional reader
that distinguishes them: a column sum has fewer coordinates
than the triple that produces it. Accordingly, the temporal row used below
jointly retains the three bands and the four fields
$q,a,c,\bar w$.

\subsection{Conversion clock and first zero advance}

\subsubsection{Ternary capacity and decimal depth}

An expression of $t+1$ bands contains $6(t+1)$ trits. The associated conversion calendar is defined by the integer
\begin{equation}
 d_t=\max\{k\in\mathbb N:1000^k\leq3^{6(t+1)}\}
     =\max\{k\in\mathbb N:1000^k\leq729^{t+1}\}.
 \label{eq:kdes-reloj}
\end{equation}
The definition compares integer capacities. The counter $t$ and the
depth $d_t$ record different quantities: a transition adds
one six-trit band; the clock counts complete powers of the
decimal base one thousand contained in the accumulated ternary capacity.

\begin{lemma}[Integer characterization of the clock]
For $t,k\in\mathbb N$,
\[
 d_t=k\quad\Longleftrightarrow\quad
 1000^k\leq729^{t+1}<1000^{k+1}.
\]
Moreover, $d_t\leq d_{t+1}$.
\end{lemma}
\begin{proof}
The powers of one thousand form a strictly increasing, unbounded
sequence. The power $729^{t+1}$ lies between two consecutive terms;
the lower exponent is precisely the maximum in
\eqref{eq:kdes-reloj}. The inequality
$729^{t+1}<729^{t+2}$ implies monotonicity of the maximum.
\end{proof}

\begin{proposition}[First zero advance]
The least time satisfying $d_{t+1}=d_t$ is
\[
 t_*=20.
\]
In particular,
\[
 d_0=0,\quad d_{19}=19,\quad d_{20}=20,
 \quad d_{21}=20,\quad d_{22}=21.
\]
\end{proposition}
\begin{proof}
The exact integer comparisons
\[
 1000^{20}\leq729^{21},\qquad
 729^{22}<1000^{21},\qquad
 1000^{21}\leq729^{23}<1000^{22}
\]
give the values at times twenty, twenty-one, and twenty-two. For
$0\leq t\leq20$, the expression
$729^{t+1}/1000^t=729(729/1000)^t$ is decreasing and its value at twenty
is greater than or equal to one. Hence,
$1000^t\leq729^{t+1}<1000^{t+1}$ and $d_t=t$. No time less than
twenty can satisfy $d_{t+1}=d_t$; at twenty it does hold.
All comparisons used involve integer powers.
\end{proof}

The number twenty is therefore an output of the minimality criterion. The rule
selecting the instant is ``the first zero advance of the clock''; the
algorithm and its proof can use this rule without receiving twenty as
an external parameter.

\subsubsection{Nonadic return and temporal antecedent}

The length of the phase return determines
\begin{equation}
 r=t_*-9=11.
 \label{eq:kdes-retorno}
\end{equation}
The descriptor queries are made at $t_*+1$, $r+1$, and $r$.
This distribution preserves the relation between the pause instant, the
nonadic translation, and the immediately associated phase reading. Times
eleven, twelve, and twenty-one designate positions in the regional
calendar; their temporal difference remains recorded even when two
times produce the same decimal depth.

\subsection{Oriented construction of the twenty-four-trit descriptor}

\subsubsection{Opposite half-band defects}

Take the oriented defects of the regional chart
\[
 \delta_+=(0,0,0,1,1,1),\qquad
 \delta_-=(0,0,0,2,2,2)\in\mathbb F_3^6.
\]
Their opposition is verified componentwise:
\[
 \delta_++\delta_-=0\quad\text{in }\mathbb F_3^6.
\]
The representative $2$ realizes $-1$ in the ternary chart. Concatenating the three-position tails in negative--positive order produces the word $222111$. Summing defects and concatenating their tails are distinct operations: the former has codomain $\mathbb F_3^6$, whereas the latter preserves the order of two segments of length three.

\subsubsection{Dual charge and phase selection}

In the regional chart, the six-coordinate transformation uses
\[
 A_{\rm reg}=
 \begin{pmatrix}
 0&1&1&1&1&1\\
 1&0&1&1&2&2\\
 1&1&0&2&1&2\\
 1&1&2&0&2&1\\
 1&2&1&2&0&1\\
 1&2&2&1&1&0
 \end{pmatrix}.
\]
With row vectors, the dual charge of $b\in\mathbb F_3^6$ is the sum
of the coordinates of $bA_{\rm reg}$. The integer row sums of
$A_{\rm reg}$ are $(5,7,7,7,7,7)$; this gives the local expression
\begin{equation}
 \eta^\vee(b)=2b_0+b_1+b_2+b_3+b_4+b_5\pmod3.
 \label{eq:kdes-cargadual}
\end{equation}
The reading phase at instant $r$ is
$\theta_r=(r-1)\bmod3$. For each channel, define
\[
 \varepsilon_\chi(r)=
 \begin{cases}
 1,&\eta^\vee(B_\chi(r))=\theta_r,\\
 2,&\eta^\vee(B_\chi(r))\ne\theta_r.
 \end{cases}
\]
The phase--charge expression is
\begin{equation}
 D(r)=\varepsilon_\pi(r)B_\pi(r)
      +\varepsilon_e(r)B_e(r)
      +\varepsilon_\varphi(r)B_\varphi(r)
      \quad\text{in }\mathbb F_3^6.
 \label{eq:kdes-fasecarga}
\end{equation}
The coefficient is determined by an equality of charges before evaluating
the word that will result from the combination.

\subsubsection{Composition of the descriptor}

Let $\Vert$ denote the concatenation of fixed-width words.
The complete construction is
\begin{equation}
 \begin{split}
 \Omega_{24}={}&
 (B_e(t_*+1)+\delta_-)
 \Vert (\delta_-|_{3,4,5}\Vert\delta_+|_{3,4,5})\\
 &\Vert(B_\pi(r+1)+B_e(r+1))\Vert D(r).
 \end{split}
 \label{eq:kdes-omega}
\end{equation}
Each length-six summand is evaluated in $\mathbb F_3^6$; the
final concatenation has length $6+6+6+6=24$. This construction
distinguishes the descriptor $\Omega_{24}$ from the regional stages
$w_{24}^{\chi}$: the two objects share a length, but their producers
and their places in the composition differ.

\begin{proposition}[Regional evaluation of the descriptor]
The regional expressions and the first-zero-advance criterion determine
\[
 \Omega_{24}=020022\,222111\,211201\,021101,
 \qquad [\Omega_{24}]_3=66241910521.
\]
\end{proposition}
\begin{proof}
The required bands, preserving the order $(\pi,e,\varphi)$, are
\[
 \begin{array}{c|ccc}
 t&B_\pi(t)&B_e(t)&B_\varphi(t)\\\hline
 11&022212&110001&100021\\
 12&012012&202222&000120\\
 21&221022&020100&211021
 \end{array}
\]
and come from extraction \eqref{eq:kdes-extraccion}. At time
twenty-one,
$020100+000222=020022$ in $\mathbb F_3^6$. The defect tails
provide $222111$. At time twelve,
$012012+202222=211201$.

At time eleven, the dual charges are $0,1,2$, respectively;
the phase is $(11-1)\bmod3=1$. The resulting coefficients are $(2,1,2)$ and
\[
 2(022212)+(110001)+2(100021)=021101
 \quad\text{in }\mathbb F_3^6.
\]
Concatenating these four results gives the stated word.
Its integer evaluation is obtained by Horner's method in base three, or by three
successive concatenations in base $729$.
\end{proof}

The preceding table serves a verification purpose: each of its
entries is obtained from the regional producer. Construction rule
\eqref{eq:kdes-omega} makes it possible to reproduce the descriptor without
introducing the final twenty-four-trit string as an input.

\subsection{Functional panel and terminal repertoire}

\subsubsection{Nine regional positions}

The functional panel retains the first three bands of each of the
three channels, in the declared regional order:
\begin{equation}
 \mathcal P_9=
 \bigl(b_\pi(0),b_\pi(1),b_\pi(2),
       b_e(0),b_e(1),b_e(2),
       b_\varphi(0),b_\varphi(1),b_\varphi(2)\bigr).
 \label{eq:kdes-panel}
\end{equation}
Its evaluation is
\[
 \mathcal P_9=(103,161,103,523,457,301,450,398,438).
\]
The nine positions retain their regional provenance, although the value
103 occurs twice. The panel has nine ordered incidences and
eight distinct values. The subsequent enumeration uses the incidences;
candidate deduplication takes place at a later stage.

\subsubsection{Unordered eight-block repertoire}

The terminal repertoire involved in this construction is
\begin{equation}
 \mathcal S_{\rm term}=
 \{002001,020111,200110,121012,
   010122,001001,221110,222110\}\subset\mathbb F_3^6.
 \label{eq:kdes-sterm}
\end{equation}
Its eight elements are distinct. The integer encoding, sorted by
value solely to fix a computational representation, is
\[
 \{28,55,98,175,437,498,687,714\}.
\]
The documentary provenance of this repertoire is the terminal segmentation
of the seventy-two-trit word and its position-wise recovery
through incidences of the phase--charge and affine readers. The selector
developed below receives the repertoire without its historical
ordering and tests its $8!$ orderings.

The exact formulation of the reassembled implementation maintains an
explicit interface here: the regional producers materially replace
the descriptor, the panel, and the temporal rows; repertoire
\eqref{eq:kdes-sterm} retains its earlier documentary producer. This
distinction identifies where each input enters. Membership of
eight words in a temporal reader repertoire, selection of those
eight words, and determination of their terminal order are three
different statements. The terminal-selection proof establishes the
third for the declared repertoire.

Documentary recovery by position can be expressed
exactly. Let $\mathcal I$ be the phase--charge incidence register and let
$\mathcal J$ be the register of the affine extension. Each incidence retains
the position $\ell$ in the seventy-two-trit word and the
emitted word $u$. For $5\leq\ell\leq12$, form
\[
 \mathcal S_\ell=
 \{u:(\ell,u)\in\mathcal I\cup\mathcal J\}.
\]
The recovery procedure checks that all eight positions occur and that each
$\mathcal S_\ell$ is a singleton. If $s_\ell$ is its unique element,
the output delivered to the selector is
\[
 \mathcal S_{\rm term}=\{s_5,s_6,\ldots,s_{12}\}.
\]
The singleton proof for each position preserves the coherence of the
historical witnesses; the subsequent export retains the eight
words and removes only the historical ordering. The selector
reconstructs the order through complete enumeration and the terminal
conditions, instead of receiving it from the provenance register.

% RESULTADO_RECUPERADO: K31-10/P10 del inventario REV05.
% Fuente de las incidencias:
% 16_CIERRE_GLOBAL_HMT_MD_2026-07-22/02_LECTURA_DODECAFASICA/continuacion_fases/
% verificar_obstruccion_y_dato_minimo_w72.py, lineas 45--60 y 141--233.
% RESULTADO_OBSTRUCCION_Y_DATO_MINIMO_W72.json, phase_charge_reader.
% Fuente de las bandas: N69_signatures_t0_t99.csv, tiempos impresos abajo.
% Consumidor: PAQUETE_CONTINUIDAD_K_UNIDAD_20260918_BASE_COMPARTIDA/python/
% selector_orbital_original.py, verify_historical_inputs, lineas 159--193.
% No se transfiere la restriccion del lector comparativo al generador completo.
\subsubsection{Temporal incidences of the terminal repertoire and recovery by position}
\label{sec:repertorio-incidencias-completas}

The expression of the terminal repertoire uses the same regional blocks, carries and calendar that produce the initial twenty-four-trit descriptor. Two operations in this composition should be made explicit separately. The first recovers eight words from recorded temporal incidences. The second presents those eight words to the dodecaphase selector as an unordered set. The selector re-establishes their order through its incidence, cylinder and closure conditions. This distinction preserves provenance witnesses without turning the order of a historical segmentation into an input to the selector.

Let $B_t=(b_{t,0},b_{t,1},b_{t,2})\in(\mathbb F_3^6)^3$ be the regional boundary already generated at time $t$, with $0\leq t<100$. In the chart used by the historical register, the incidence matrix is
\[
 A_{\mathrm{hist}}=
 \begin{pmatrix}
 0&1&1&1&1&1\\
 1&0&1&1&2&2\\
 1&1&0&2&1&2\\
 1&1&2&0&2&1\\
 1&2&1&2&0&1\\
 1&2&2&1&1&0
 \end{pmatrix}.
\]
The words are regarded as row vectors. Their two charges and the
chronological phase are
\[
 r_{t,i}=\sum_{j=1}^{6}b_{t,i,j},\qquad
 d_{t,i}=\sum_{j=1}^{6}(b_{t,i}A_{\mathrm{hist}})_j,
 \qquad \vartheta_t=t-1\pmod 3.
\]
All operations in these three expressions take place in $\mathbb F_3$. The historical subscript identifies a concrete chart of Witt incidence; it keeps the coordinate transformation explicit when another presentation of that same incidence is used.

For $c\in\{r_t,d_t\}$, $\delta\in\mathbb F_3$ and
$\epsilon\in\{1,-1\}$, define
\[
 \eta_i(c,\delta,\epsilon)=
 \begin{cases}
  \epsilon,&c_i=\vartheta_t+\delta,\\
  -\epsilon,&c_i\ne\vartheta_t+\delta,
 \end{cases}
 \qquad
 L_{t,c,\delta,\epsilon}
   =\sum_{i=0}^{2}\eta_i(c,\delta,\epsilon)b_{t,i}.
\]
The orientation $-1$ is written as $2$ when printing the ternary
coefficients. Alongside these twelve comparative readers per boundary, the
register has the affine reader
\[
 L^{\mathrm{af}}_t
    =\sum_{i=0}^{2}(d_{t,i}+\vartheta_t)b_{t,i}.
\]
Adding the phase in the affine reader and comparing charges in
the preceding reader are different operations, both acting on data already
produced by the TPK boundary.

\begin{table}[htbp]
\centering\small
\begin{tabular}{rccc cc}
\toprule
$t$&$b_{t,0}$&$b_{t,1}$&$b_{t,2}$&$r_t$&$d_t$\\
\midrule
8 &200121&212212&110110&011&202\\
11&022212&110001&100021&001&012\\
30&210112&110202&211110&100&012\\
34&121022&000011&112110&220&021\\
37&120121&011202&112221&100&201\\
49&110212&200212&100122&110&201\\
58&111121&022121&122201&122&220\\
67&010012&001220&011120&122&122\\
75&220000&020111&100002&120&021\\
81&122201&020202&201122&202&001\\
90&022112&022110&121010&202&200\\
\bottomrule
\end{tabular}
\caption{Temporal boundaries realizing the printed incidences of the
terminal repertoire. Each six-digit entry preserves the order of
the six coordinates of the regional boundary.}
\label{tab:terminal-bandas-testigo}
\end{table}

The following table displays all comparative incidences of the
register on the twelve blocks of the segmentation under consideration. The
letters $r$ and $d$ denote the visible charge and the dual charge,
respectively. Column $j$ retains the position index of the historical witness;
the operation delivering the data to the selector forgets this index after forming
the set of words.

\begin{table}[htbp]
\centering\small
\begin{tabular}{rrcrrc c}
\toprule
$j$&$t$&charge&$\delta$&$\epsilon$&$(\eta_0\eta_1\eta_2)$&$L_{t,c,\delta,\epsilon}$\\
\midrule
4&11&$d$&0& 1&212&021101\\
5&67&$d$&1&$-1$&211&002001\\
5&67&$d$&2& 1&211&002001\\
5&67&$r$&1&$-1$&211&002001\\
5&67&$r$&2& 1&211&002001\\
6&81&$d$&0& 1&222&020111\\
6&81&$r$&2& 1&222&020111\\
7&34&$r$&1&$-1$&111&200110\\
7&75&$d$&1&$-1$&211&200110\\
7&75&$r$&2&$-1$&211&200110\\
8&90&$r$&0& 1&121&121012\\
8&90&$r$&1&$-1$&121&121012\\
9&49&$d$&0&$-1$&121&010122\\
11&11&$d$&2&$-1$&211&221110\\
11&37&$d$&0&$-1$&121&221110\\
12&8&$d$&0&$-1$&111&222110\\
12&8&$r$&1&$-1$&111&222110\\
12&58&$d$&1&$-1$&111&222110\\
12&58&$r$&0&$-1$&111&222110\\
\bottomrule
\end{tabular}
\caption{Complete phase--charge incidences on the terminal segmentation
within the hundred boundaries of the register. Word matches between
rows retain their times, orientations, and charge types.}
\label{tab:terminal-incidencias-completas}
\end{table}

The tenth block has an affine witness. At $t=30$,
$d_{30}=(0,1,2)$ and $\vartheta_{30}=2$, whence
\[
 (d_{30,0}+\vartheta_{30},d_{30,1}+\vartheta_{30},
 d_{30,2}+\vartheta_{30})=(2,0,1)
\]
and, using the bands in Table~\ref{tab:terminal-bandas-testigo},
\[
 L^{\mathrm{af}}_{30}
 =2(210112)+0(110202)+(211110)
 =(001001)\quad\text{in }\mathbb F_3^6.
\]
The equality holds coordinatewise: before reduction modulo three,
the sum vector is $(6,3,1,3,3,4)$.

\begin{proposition}[Complete recovery of the repertoire through incidences]
\label{prop:repertorio-terminal-ocho-incidencias}
Consider the comparative-incidence register in
Table~\ref{tab:terminal-incidencias-completas}, augmented by the affine
witness $(j,t,L)=(10,30,001001)$. For each position $5\leq j\leq12$, the
set of output words of its incidences is a singleton.
Extracting these eight words and subsequently forgetting their order
produces exactly
\[
 \mathcal S_{\mathrm{term}}=
 \{002001,020111,200110,121012,010122,001001,221110,222110\}.
\]
Their encoding as base-three integers is
\[
 \{28,55,98,175,437,498,687,714\}.
\]
\end{proposition}
\begin{proof}
The incidence multiplicities at positions $j=5,6,7,8,9,10,11,12$
are, respectively,
\[
 (4,2,3,2,1,1,2,4).
\]
At each position, all rows print the same word. Each row
is verified by substituting its three bands and three coefficients into the
formula for $L$; the tenth case has been calculated explicitly through
$L^{\mathrm{af}}_{30}$. For example, at $t=67$ and with coefficients
$(2,1,1)$, one obtains
\[
 2(010012)+(001220)+(011120)=(002001)\pmod3.
\]
The eight results are distinct as words of length six. Their
encoding $\sum_{k=1}^{6}w_k3^{6-k}$ gives, in position
order, $(55,175,498,437,98,28,714,687)$. Forgetting the order
yields the stated set of integers. Thus, extraction preserves
each complete word, including its zero positions, and delivers eight
distinct elements to the selector.
\end{proof}

The preceding proof is an exact recovery from identified
incidences. Its input comprises the incidence register
with its positions and witnesses. The subsequent orbital selection receives
only $\mathcal S_{\mathrm{term}}$ and the initial regional descriptor;
it traverses the $8!$ orderings and determines their compatibility with the
functional boundaries and the oriented incidences of the dodecaphasic
wheel. Both operations are thus presented contiguously, with explicit
domains: temporal provenance of the eight words and constructive
determination of their order.

The comparative table contains the nineteen incidences that its formula
produces on the twelve-block segmentation over the hundred declared
times. This finite exhaustiveness is checked by traversing
$100\cdot2\cdot3\cdot2=1200$ evaluations. The affine reader extends
coverage of the terminal positions from seven to eight. This count
characterizes the local reader repertoire just defined;
the other TPK operations retain the domains and constructive functions
established in their respective derivations.

\subsection{Comparative readers, affine reader, and temporal incidence}

\subsubsection{Witt transformation of the selector}

The selector uses the six-coordinate chart
\[
 A_{\rm ter}=
 \begin{pmatrix}
 0&1&1&1&1&1\\
 1&0&1&2&2&1\\
 1&1&0&1&2&2\\
 1&2&1&0&1&2\\
 1&2&2&1&0&1\\
 1&1&2&2&1&0
 \end{pmatrix}.
\]
Define $W(b)=bA_{\rm ter}$ in $\mathbb F_3^6$, together with the charges
\[
 \eta(b)=\sum_{j=0}^5 b_j\pmod3,
 \qquad\eta_{\rm ter}^\vee(b)=\eta(W(b)).
\]
The descriptor chart and selector chart are related by the permutation $\sigma=(0,1,2,4,5,3)$, written as a list of images:
\begin{equation}
 (A_{\rm ter})_{ij}=(A_{\rm reg})_{\sigma(i),\sigma(j)}.
 \label{eq:kdes-cambio-carta}
\end{equation}

\begin{lemma}[Compatibility of the charge functional]
For every band $b$,
\[
 \eta_{\rm ter}^\vee(b)=\eta^\vee(b).
\]
\end{lemma}
\begin{proof}
The integer sum of each corresponding row of the two matrices agrees:
it is five in the first row and seven in the other five. By
distributivity,
\[
 \sum_j\sum_i b_i(A_{\rm ter})_{ij}
 =\sum_i b_i\sum_j(A_{\rm ter})_{ij}
 =\sum_i b_i\sum_j(A_{\rm reg})_{ij}.
\]
Reduction modulo three gives the equality of charges. Matrix identity
\eqref{eq:kdes-cambio-carta} is verified by substituting its
thirty-six entries and also preserves the coordinate change.
\end{proof}

Equality of the total charge functional allows their expressions to be compared. Vector transformations remain identified by their matrices and by the explicit permutation; the preceding proof preserves that coordinate information.

\subsubsection{Twelve comparative readers per instant}

Define the temporal phase
\[
 \theta(t)=(t+2)\bmod3.
\]
This expression realizes $(t-1)\bmod3$ also at $t=0$, where the phase
equals two. For each charge type
$\nu\in\{\eta,\eta_{\rm ter}^\vee\}$, offset
$o\in\{0,1,2\}$ and orientation $s\in\{1,2\}$, take
\[
 \epsilon_\chi^{\nu,o,s}(t)=s
 \begin{cases}
 1,&\nu(B_\chi(t))=\theta(t)+o\pmod3,\\
 2,&\nu(B_\chi(t))\ne\theta(t)+o\pmod3.
 \end{cases}
\]
The corresponding comparative word is
\begin{equation}
 C_{\nu,o,s}(t)=
 \sum_{\chi\in\{\pi,e,\varphi\}}
 \epsilon_\chi^{\nu,o,s}(t)B_\chi(t)
 \quad\text{in }\mathbb F_3^6.
 \label{eq:kdes-comparativo}
\end{equation}
The parameters generate twelve row-labelled expressions. Two labels may produce the same word; temporal witnesses and labels remain available in the full calendar.

\subsubsection{Affine phase--charge reader}

The affine reader uses coefficients
\[
 a_\chi^{\rm af}(t)=\eta_{\rm ter}^\vee(B_\chi(t))+\theta(t)
 \pmod3,
\]
and produces
\begin{equation}
 F(t)=\sum_{\chi\in\{\pi,e,\varphi\}}
 a_\chi^{\rm af}(t)B_\chi(t)
 \quad\text{in }\mathbb F_3^6.
 \label{eq:kdes-afin}
\end{equation}
The comparative rule distinguishes agreement and difference of charges;
the affine rule adds charge and phase before combining the bands. This operational
difference is what heterotypic selection uses.

\subsubsection{Word--residue incidence}

Let $R$ be the cyclic rotation of the six positions of a band. In the
row at time $t$, form the residual repertoire
\[
 \mathcal R(t)=\{R^jv:
      v\in\{q(t),a(t),c(t),\bar w(t)\},\ 0\leq j<6\}.
\]
The temporal incidence relations are
\begin{align}
 \mathcal A_{\rm cmp}
 &=\{(C_{\nu,o,s}(t),u):
      0\leq t<100,\ u\in\mathcal R(t),\ \nu,o,s\text{ admissible}\},\\
 \mathcal A_{\rm af}
 &=\{(F(t),u):0\leq t<100,\ u\in\mathcal R(t)\}.
 \label{eq:kdes-atlas}
\end{align}
Each pair relates a read word to a residue of the same row.
Time is quantified jointly across both coordinates of the pair.
In the calendar used, times $0,\ldots,99$ are distinct;
identifying the same row and identifying the same time are therefore
equivalent operations.

\begin{proposition}[Existential compression of the temporal repertoire]
The function
\[
 M(v,u)=\mathbf1_{(v,u)\in\mathcal A_{\rm cmp}}
        +2\mathbf1_{(v,u)\in\mathcal A_{\rm af}}
        \in\{0,1,2,3\}
\]
preserves exactly the existence of each reading type for the pair
$(v,u)$.
\end{proposition}
\begin{proof}
The first bit equals one exactly when there exist a time, a residual
field, a rotation, and a comparative label that produce the pair.
The second bit expresses the corresponding condition for the affine reader.
Union by bitwise disjunction preserves each existence condition
as all rows are traversed. A pair may admit both types, in which
case the mask equals three.
\end{proof}

The function $M$ is an existence reader. The witness times, their
multiplicities, and the individual labels remain in
\eqref{eq:kdes-atlas} and in the calendar that produces them. The
finite implementation uses $M$ to decide heterotypy, whereas
the genealogical archive retains the richer provenance data.

\subsection{Enumeration of cylinders and decimal candidates}

\subsubsection{Concatenation of seventy-two and seventy-eight trits}

For an ordering $\tau=(s_1,\ldots,s_8)$ of
$\mathcal S_{\rm term}$, define
\[
 P_{72}(\tau)=[\Omega_{24}\Vert s_1\Vert\cdots\Vert s_8]_3.
\]
Recursively, start from $p_0=[\Omega_{24}]_3$ and compute
$p_i=729p_{i-1}+[s_i]_3$ for $1\leq i\leq8$. This recurrence
preserves the leading zeros of each block through explicit
multiplication by $729$. For each position $b$ of the functional panel,
\[
 P_{78}(\tau,b)=729P_{72}(\tau)+b.
\]
The lengths are $24+8\cdot6=72$ and $72+6=78$, respectively.
The corresponding ternary cylinder is
\begin{equation}
 I_{\tau,b}=
 \left[\frac{P_{78}(\tau,b)}{3^{78}},
       \frac{P_{78}(\tau,b)+1}{3^{78}}\right).
 \label{eq:kdes-cilindro}
\end{equation}

\subsubsection{Exact intersection with the twelve-triad grid}

The decimal reading uses twelve triads, that is, the grid
$N/10^{36}$. Write $D=10^{36}$ and $Q=3^{78}$ for brevity.

\begin{lemma}[Integer endpoints of the cylinder]
For an integer prefix $p$ of length seventy-eight, the numerators
of the grid points belonging to its cylinder are
\begin{equation}
 \left\{N\in\mathbb Z:
 \left\lceil\frac{pD}{Q}\right\rceil
 \leq N<
 \left\lceil\frac{(p+1)D}{Q}\right\rceil\right\}.
 \label{eq:kdes-rejilla}
\end{equation}
Each cylinder contains at most one point of this grid.
\end{lemma}
\begin{proof}
The condition $N/D\in[p/Q,(p+1)/Q)$ is equivalent to
$pD/Q\leq N<(p+1)D/Q$. For integer $N$, the first inequality
is equivalent to $\lceil pD/Q\rceil\leq N$; the second is equivalent to
$N<\lceil(p+1)D/Q\rceil$. The interval length in coordinate
$N$ is $D/Q<1$, since $10^{36}<3^{78}$. It can therefore contain
at most one integer.
\end{proof}

Ceilings are evaluated by integer division:
$\lceil a/Q\rceil=\lfloor(a+Q-1)/Q\rfloor$ for $a\geq0$.
The entire enumeration uses exact integers; the choice of boundary
points is fixed by the half-open right endpoint.

First retain the indexed list of incidences
\[
 \mathcal C_{\rm ind}=
 \{(\tau,j,N):N/D\in I_{\tau,\mathcal P_9(j)}\},
 \qquad1\leq j\leq9,
\]
and then obtain the set of numerators
\[
 \mathcal C=\{N:\exists\tau,j\ (\tau,j,N)\in\mathcal C_{\rm ind}\}.
\]
Passing to $\mathcal C$ removes repetitions of a numerator, retained
in $\mathcal C_{\rm ind}$ together with their provenance.

\begin{lemma}[Recovery of the ordering from the cylinder point]
If $(\tau,j,N)\in\mathcal C_{\rm ind}$, block number $i$ of the
seventy-eight-trit word is recovered by
\[
 \left\lfloor\frac{N\,729^i}{D}\right\rfloor\bmod729,
 \qquad1\leq i\leq13.
\]
In particular, blocks five through twelve recover $\tau$.
\end{lemma}
\begin{proof}
Membership in the cylinder means
$p\leq QN/D<p+1$, where $p=P_{78}(\tau,\mathcal P_9(j))$.
Hence, $\lfloor QN/D\rfloor=p$. Extraction of the
earlier blocks agrees with successive division of the prefix $p$ by
powers of $729$: truncating a positional expansion before extracting
a block preserves that block. The displayed formula performs
exactly this extraction. The first four blocks form
$\Omega_{24}$ and the next eight are the ordering $\tau$.
\end{proof}

\begin{proposition}[Exact candidate count]
For $\Omega_{24}$, $\mathcal S_{\rm term}$ and $\mathcal P_9$
defined above,
\[
 \#\operatorname{Ord}(\mathcal S_{\rm term})=40320,
 \qquad\#\mathcal C_{\rm ind}=21902,
 \qquad\#\mathcal C=19446.
\]
\end{proposition}
\begin{proof}
The eight blocks are distinct, so their orderings are the
$8!=40320$ permutations. For each permutation and each of the
nine panel positions, compute $p=P_{78}(\tau,b)$, apply
endpoints \eqref{eq:kdes-rejilla}, and add the unique integer when
the integer interval is nonempty. The sum of the cardinalities of these
intervals is $21902$. The union of their integers contains $19446$
distinct elements. The exhaustive evaluation of these operations
is certified by the theorems
\texttt{indexed\_cylinder\_count} and \texttt{candidate\_count} of the
finite selector. The algorithm performs exactly the enumeration
described: permutations, nine panel incidences, bounds by
integer division, and duplicate removal.
\end{proof}

The numbers $21902$ and $19446$ count different objects. The former
counts indexed incidences containing a decimal point; the latter
counts distinct numerators. The total number of cylinder queries
is $40320\cdot9=362880$, including cylinders whose intersection with
the grid is empty.

\subsection{Exceptional incidence and heterotypy condition}

\subsubsection{Exceptional face produced by the panel}

A band is lifted to twelve coordinates by
\[
 \Gamma(b)=(b,W(b))\in\mathbb F_3^{12}.
\]
The support of a word is the set of positions of its
nonzero coordinates, with human-readable indices $1,\ldots,12$. The first
regional bands of the panel produce
\begin{equation}
 \mathcal B=\operatorname{supp}\Gamma(B_\pi(0))
          \cap\operatorname{supp}\Gamma(B_e(0))
          =\{4,6,9,10\}.
 \label{eq:kdes-cara}
\end{equation}
The operation uses the regional words $103$ and $523$ and the
selector's Witt matrix. The face is computed before querying the
candidate coordinates.

For $N\in\mathcal C$, its twelve triads are
\begin{equation}
 K_i(N)=\left\lfloor\frac{N}{1000^{12-i}}\right\rfloor\bmod1000,
 \qquad1\leq i\leq12.
 \label{eq:kdes-triadas}
\end{equation}
The crown support is
\[
 \mathcal H(N)=\{i:K_i(N)\geq729\},
 \qquad
 \mathcal C_{\mathcal B}=\{N\in\mathcal C:\mathcal H(N)=\mathcal B\}.
\]
The boundary $729=3^6$ relates the six-trit band to the decimal
triad. Exhaustive evaluation of the equality of supports gives
\begin{equation}
 \#\mathcal C_{\mathcal B}=79.
 \label{eq:kdes-79}
\end{equation}
The support preserves positions, while the complete register
also preserves the twelve magnitudes. The next stage uses this
positional information together with relations between magnitudes.

\subsubsection{Ternary words and oriented differences of the candidate}

The $i$th ternary band of the point $N/D$ is obtained by
\begin{equation}
 T_i(N)=\left\lfloor\frac{N\,729^i}{D}\right\rfloor\bmod729,
 \qquad1\leq i\leq13.
 \label{eq:kdes-tercandidato}
\end{equation}
The oriented edge $a\to b$ has residue
\begin{equation}
 R_{a,b}(N)=(K_a(N)-K_b(N))\bmod729.
 \label{eq:kdes-resarista}
\end{equation}
Subtraction takes place in $\mathbb Z$ before reduction. For this edge, the pair is looked up in the relations \eqref{eq:kdes-atlas}
\[
 \bigl(\operatorname{enc}_6^{-1}(T_b(N)),
  \operatorname{enc}_6^{-1}(R_{a,b}(N))\bigr).
\]
It is therefore a joint query on the destination word and the oriented difference of the triads.

Write $\mathrm{Cmp}(a,b;N)$ when the pair belongs to
$\mathcal A_{\rm cmp}$ and $\mathrm{Af}(a,b;N)$ when it belongs to
$\mathcal A_{\rm af}$.

\subsubsection{Determination of the transverse orbit}

The twenty-four-trit descriptor determines the first three
decimal triads before the terminal filtrations. Indeed, its
cylinder satisfies the exact inclusion
\[
 \left[\frac{66241910521}{3^{24}},
       \frac{66241910522}{3^{24}}\right)
 \subseteq
 \left[\frac{234543140}{10^9},
       \frac{234543141}{10^9}\right).
\]
The two endpoint inequalities are verified by integer cross-products. Every candidate therefore shares the three reference positions $1,2,3$ and their expressions $(234,543,140)$.

The step-four translation on $\mathbb Z/12\mathbb Z$ has the
four orbits
\[
 (1,5,9),\qquad(2,6,10),\qquad(3,7,11),\qquad(4,8,12).
\]
The first three each contain one of the reference
positions; the fourth is the only one disjoint from $\{1,2,3\}$. The oriented
forest of steps three and four is fixed by the edges
\[
 \begin{array}{lll}
 1\to4,&1\to5,&2\to6,\\
 3\to7,&4\to8,&5\to9,\\
 6\to10,&7\to11,&8\to12.
 \end{array}
\]
Its first edge has step three and the remaining edges have step four. The two
internal edges of the distinguished orbit are precisely
$4\to8$ and $8\to12$. The location of the reading condition is thus
fixed by the geometry of the positions and by the initial prefix,
before examining the candidates of the exceptional face.

\subsubsection{Terminal heterotypy condition}

The terminal condition is
\begin{equation}
 \begin{split}
 \mathcal T(N)={}&
 [\mathrm{Cmp}(4,8;N)\wedge\mathrm{Af}(8,12;N)]\\
 &\vee[\mathrm{Af}(4,8;N)\wedge\mathrm{Cmp}(8,12;N)].
 \end{split}
 \label{eq:kdes-heterotipia}
\end{equation}
Both assignments of types to the two edges are admitted. The
selection requires the existence of a realization with distinct types; the
individual labels and their times remain available in the
incidence witnesses.

\begin{theorem}[Terminal selection of the coupling register]
On the candidate domain produced by
$\Omega_{24}$, $\mathcal S_{\rm term}$, $\mathcal P_9$ and the hundred-row
regional calendar, there exists a unique numerator simultaneously satisfying
exceptional incidence and heterotypy:
\[
 \{N\in\mathcal C_{\mathcal B}:\mathcal T(N)\}
 =\{234543140729659824621058914794146601\}.
\]
Its twelve-triad register is
\begin{equation}
 K=(234,543,140,729,659,824,621,058,914,794,146,601).
 \label{eq:kdes-kfinal}
\end{equation}
\end{theorem}
\begin{proof}
The preceding count produces the $79$ elements of
$\mathcal C_{\mathcal B}$. For each one, its twelve
coordinates are extracted by \eqref{eq:kdes-triadas}; the destination bands
$T_8,T_{12}$ and residues $R_{4,8},R_{8,12}$ are computed; the
two masks of the temporal repertoire are queried; and disjunction
\eqref{eq:kdes-heterotipia} is applied. The resulting list contains exactly
the stated integer, according to the exhaustive evaluation certified by
\texttt{terminal\_selection}. Equality of the list with a singleton
implies existence and uniqueness in the declared domain. Base-one-thousand
extraction from that integer gives \eqref{eq:kdes-kfinal}.

The selection procedure receives the enumerated producers and rules.
The final integer appears in the statement of the verified
equality after selection has been executed. This logical position
allows the result to be checked without using it as a selection criterion.
\end{proof}

\subsubsection{Arithmetic witness of membership in a cylinder}

One of the incidences of the selected numerator uses the ordering
\[
 (002001,020111,200110,121012,010122,001001,221110,222110)
\]
and the boundary $438=121020_3$. For this incidence,
\[
 P_{72}=5283881584882673136514102926090957,
\]
\[
 P_{78}=3851949675379468716518781033120308091.
\]
The integer-endpoint formula yields
\begin{align*}
 \left\lceil\frac{P_{78}10^{36}}{3^{78}}\right\rceil
 &=234543140729659824621058914794146601,\\
 \left\lceil\frac{(P_{78}+1)10^{36}}{3^{78}}\right\rceil
 &=234543140729659824621058914794146602.
\end{align*}
The unique integer in the interval is therefore $N_K$. The thirteen ternary
bands of its decimal point are
\[
 \begin{array}{rrrr}
 020022&222111&211201&021101\\
 002001&020111&200110&121012\\
 010122&001001&221110&222110\\
 121020&&&
 \end{array}
\]
This calculation exhibits a concrete witness of membership in the enumerated
domain. Uniqueness of the register follows from the exhaustive examination of
all candidates, in addition to this particular witness.

\subsection{Regional composition and continuity of the realizations}

\subsubsection{Composition of the regional generators}

Let $\operatorname{Sel}(\Omega,\mathcal S,\mathcal P,\mathcal N)$ be the
selector defined by the enumeration and the two preceding filters.
The regional producers satisfy the equalities
\[
 \Omega_{\rm reg}=\Omega_{24},\qquad
 \mathcal P_{\rm reg}=\mathcal P_9,\qquad
 \mathcal N_{\rm reg}=\mathcal N_{100},
\]
where $\mathcal N_{\rm reg}$ is computed by band extraction and
signature reading. By congruence of functions,
\begin{equation}
 \begin{split}
 &\operatorname{Sel}(\Omega_{\rm reg},\mathcal S_{\rm term},
                    \mathcal P_{\rm reg},\mathcal N_{\rm reg})\\
 &\hspace{10mm}=
 \operatorname{Sel}(\Omega_{24},\mathcal S_{\rm term},
                    \mathcal P_9,\mathcal N_{100})=\{N_K\}.
 \end{split}
 \label{eq:kdes-composicion}
\end{equation}
The equality replaces three documentary interfaces with their actual
regional construction. The terminal repertoire retains the provenance
specified in \eqref{eq:kdes-sterm}.

The register is then formed by the twelve Euclidean divisions in
\eqref{eq:kdes-triadas}. Each coordinate belongs to
$\{0,\ldots,999\}$ and the length is twelve by construction. The
implementation that takes the first element of the selected list
is justified by equality with the singleton: the default value
of a computational operation on empty lists never enters
this result.

\subsubsection{Relation to signed incidences and inversion}

The present construction provides a prospective selection in an
explicit finite domain. The subsequent development of the dodecaphasic
register produces its oriented channels from the enriched terminal
state and proves an integral inversion through
Hadamard. That inversion recovers a register from compatible
channels; here a register is selected from its
regional and incidence-based inputs. Equality of their results
allows the two realizations to be composed while preserving the origin of each
operation.

Uniqueness of selection and uniqueness of inversion are distinct
properties. The former quantifies over the candidates constructed in this
chapter. The latter quantifies over the registers sharing a compatible
signature. Their conjunction justifies transporting the selected register
to the signed realization and recovering it from that realization, without
turning a retrospective recovery into the producer of the register.

\subsubsection{Relation between the terminal counts}

The sequence
\[
 40320\ \longrightarrow\ 21902\ \longrightarrow\ 19446
 \ \longrightarrow\ 79\ \longrightarrow\ 1
\]
records, successively, orderings, indexed incidences with the
grid, distinct numerators, candidates of the exceptional face, and
heterotypic survivors. The arrow notation indicates the order of
the operations; the second number counts a different class of objects
from the first.

The regional-catalogue census $196\cdot197\cdot196\longrightarrow331\longrightarrow2
\longrightarrow1$, developed in its corresponding place, starts from a different population: regional blocks from three catalogues and terminal Gram, occupation, distance and orientation conditions. Both censuses preserve their domains. Their convergence in a common expression is conveyed through identification of their readers and outputs; intermediate figures describe different operations and remain documented separately.

\subsubsection{Regional composition of the fine-structure constant}

The output $K$ then enters the ordered regional composition and the base-thousand carry. Its subsequent use preserves the twelve positions and the reading boundary. Coupling to the analytic root and expressions at arbitrary depth uses that same output, together with the regional coordinates already generated and the rules of the analytic chart. This dependency prepares the chapter devoted to the two correlated realizations of the fine-structure constant.

The scope of this chapter is precise: production of the descriptor and
its regional inputs, determination of the finite domain,
unique selection on that domain, and material composition with the
producers. The analytic recurrence at all orders and the
compatibility of infinite carries are developed subsequently with their
own proofs.

\subsection{Executable certification and traceability of the operations}

The formal development preserves the correspondence between the
mathematical definitions and their executable realizations. The signature
and clock proofs are found in
\texttt{N69RegionalSignature.lean}; band extraction and
reproduction of the hundred rows, in
\texttt{GeneratedN69Rows.lean}; the minimality criterion and construction
of $\Omega_{24}$, in \texttt{RegionalW24.lean}. Compatibility of
the two Witt matrices is assembled in
\texttt{WittReaderCompatibility.lean}.

The temporal readers, rotations, oriented residues, and
heterotypic condition are formalized in
\texttt{TerminalReaders.lean}. Enumeration and counts are
formalized in \texttt{Terminal}\allowbreak\texttt{Orbital}\allowbreak\texttt{Selection.lean}. The regional
panel is linked through \texttt{Regional}\allowbreak\texttt{Panel}\allowbreak\texttt{Bridge.lean}, and
composition \eqref{eq:kdes-composicion} is expressed in
\texttt{Selected}\allowbreak\texttt{From}\allowbreak%
\texttt{Regional}\allowbreak\texttt{Inputs.lean}.

The finite enumeration checks use Lean's native
evaluation; their reports record \texttt{Lean.ofReduceBool}, in addition
to the library's customary logical axioms. The regional reconstruction
of the descriptor uses kernel proofs and ordinary
reduction. This distinction describes the verification mechanisms
actually employed. The compilation receipts attest to their declared modules
and dependencies; the mathematical proofs retain the
domains, producers, and hypotheses made explicit in the body
of the chapter.
